\documentclass[10pt,a4paper]{article}
\usepackage[margin=1.55cm]{geometry}
\usepackage{xcolor}
\usepackage{amsmath,amssymb}
\usepackage{graphicx}
\usepackage{enumitem}
\usepackage{fontspec}
\usepackage[CJKmath=true]{xeCJK}
% CJK main font is resolved at COMPILE time on the actual machine, not baked in at
% generation time, so a .tex generated on one OS still renders on another (the older
% Python-side probe hardcoded one OS's font name into the file and broke when moved).
% Walk a cross-platform priority list and use the first font that exists; AutoFakeBold
% + AutoFakeSlant give bold/italic CJK even on faces that ship no bold/italic cut.
\newcommand{\setCJKto}[1]{\setCJKmainfont{#1}[AutoFakeBold=true,AutoFakeSlant=0.2]}
\IfFontExistsTF{Noto Sans CJK SC}{\setCJKto{Noto Sans CJK SC}}{%        % Linux (fonts-noto-cjk); cross-platform if installed
 \IfFontExistsTF{Source Han Sans SC}{\setCJKto{Source Han Sans SC}}{%   % Adobe 思源黑体; cross-platform
  \IfFontExistsTF{PingFang SC}{\setCJKto{PingFang SC}}{%               % macOS 10.11+
   \IfFontExistsTF{Hiragino Sans GB}{\setCJKto{Hiragino Sans GB}}{%    % older macOS
    \IfFontExistsTF{Microsoft YaHei}{\setCJKto{Microsoft YaHei}}{%     % Windows
     \IfFontExistsTF{WenQuanYi Micro Hei}{\setCJKto{WenQuanYi Micro Hei}}{% % older Linux
      \setCJKto{Songti SC}}}}}}}                                       % macOS bottom fallback
\definecolor{cblue}{HTML}{4C72B0}
\setlength{\parskip}{3pt}\setlength{\parindent}{0pt}
\setlist[enumerate]{topsep=2pt,itemsep=2pt,parsep=1pt,partopsep=0pt}
\setlist[itemize]{topsep=2pt,itemsep=2pt,parsep=1pt,partopsep=0pt}
% Let prose-embedded inline math break across lines and absorb small overflows, so simple
% formulas inserted mid-sentence stop running past the right margin.
\tolerance=1200
\emergencystretch=2.5em
\relpenalty=20
\binoppenalty=40
\hfuzz=1pt
\begin{document}

\section*{Kinetics-Clock PINNs for Endogenous High-Frequency VO2 Switching Fronts}
\textbf{Method.} Kinetics-Clock Physics-Informed Neural Network (KC-PINN)
\section*{Motivation}
The structural problem is that an endogenous switching law can compress most phase change into a short part of physical time that also moves through space. In a two-dimensional VO2 device, different locations therefore pass through the transition at different rates while voltage, temperature, and phase remain tightly coupled. A network written only in global physical time must represent a sharp front and a smooth background with the same coordinate, exposing the high-frequency and residual-conditioning weaknesses highlighted by arxiv:2602.19265v1 and the adaptive-spectral diagnosis in arxiv:2605.04502v1.

The closest architectural precedent, NeuSA (arxiv:2509.04966v2), adds a global spectral state and causal evolution for canonical wave problems, but it does not construct a spatially local clock from material kinetics. The energy-dissipation phase-field physics-informed neural network (PINN) in openalex:W4403555504 avoids propagation failure by solving a static steady-state problem, so it cannot represent pulse-created event time or moving topology. Kinetics-Clock PINN (KC-PINN) instead uses a finite analytic Gaussian basis with error-function primitives to construct the local clock without numerical time quadrature. The original electro-thermal-phase equations remain authoritative, and every first- and second-order derivative is transformed back through the learned map, including all spatial Hessian and cross terms created by a clock that varies across the device.

Three recent strands make this narrow gap actionable at once. arxiv:2605.04502v1 isolates gradient-scaling effects in adaptive spectral PINNs, arxiv:2602.19265v1 consolidates the high-frequency failure mechanism, and semanticscholar:a1448fd2600739b73894f3b8d467bf9a9e9630b2 exposes sub-nanosecond VO2 electro-thermo-phase topology as the relevant physical regime. A named FEniCSx model can supply controlled synthetic trajectories, while automatic differentiation can evaluate the complete local-coordinate transformation rather than approximate or drop its cross terms. The opportunity is therefore not merely to add spectral capacity, but to test whether an endogenous kinetics-aligned coordinate removes this specific representation bottleneck under a retrieval-bounded novelty claim.

NeuSA represents canonical initial-value problems with a truncated global spectral state evolved by a neural ordinary differential equation and causal physics constraints. It does not learn a spatially local monotone clock tied to the material phase-rate law or transform a coupled electro-thermal-phase device partial differential equation (PDE) through that clock. Its global spectral evolution coordinate lacks the endogenous, spatially heterogeneous kinetics signal needed to define local material time. The adaptive-spectral study changes Fourier frequencies and initial-condition gates to diagnose gradient scaling in a stiff nonlinear ordinary differential equation, but it does not transform time in a two-dimensional coupled device PDE or retain chain-rule terms induced by a spatially varying map. Its diagnostic setting has no spatially resolved phase-kinetic right-hand side from which an endogenous local clock could be constructed and tested. The energy-embedded phase-field PINN predicts steady ferroelectric microstructures by enforcing energy dissipation in a static formulation. It addresses dynamic PINN propagation failure by replacing the trajectory problem with a steady variational target, rather than retaining pulse-driven evolution, switching time, and moving phase-front topology. The operando VO2 study resolves electrically steered conduction topologies and sub-nanosecond electro-thermo-phase dynamics. Its contribution is physical measurement and mechanism discrimination, not a differentiable PDE solver or a training coordinate derived from measured kinetics, so it does not provide a PINN coordinate mechanism for solving the governing device fields or isolate kinetics-specific representation value on synthetic numerical evidence.

Closing the gap would make pulse-created event time, moving front topology, and electro-thermal energy closure jointly representable by one two-dimensional PINN without replacing the underlying equations with a static surrogate. This would connect the dynamic target regime exposed by semanticscholar:a1448fd2600739b73894f3b8d467bf9a9e9630b2 to a solver that can be evaluated on whole held-out devices or protocols, while directly testing whether the global spectral remedy of arxiv:2509.04966v2 is insufficient when kinetics vary across space. It would also establish a reusable conditional statement: when local phase-rate magnitude is informative and the transformation is complete, coordinate adaptation can be attributed to material kinetics rather than generic extra capacity.
\section*{Method}
\subsection*{M0\_background}
\emph{Freeze the two-dimensional electro-thermal-phase equations, pulse contract, conditions, and residual scales used by every arm.}
\begin{enumerate}[leftmargin=*,start=1]
  \item Convert the existing two-dimensional device model into a read-only `PhysicalContract`. Store its mesh and labels, material regions, pulse segments and left and right values at jumps, material functions with their inputs, starting state, boundary and interface conditions, units, and both directions of every nondimensional conversion. Make every later step read the same version. 【author decision: designate the authoritative FEniCSx model snapshot and approve its complete physical-contract manifest, reference quantities, and residual-scale rules.】 Derive the dimensionless electrical, heat, and phase operators without adding a governing term. From training cases only, create and freeze a named scale for each interior, starting, boundary, material, and pulse-interface residual. Output the operators, pulse-interface table, unit conversions, and residual-scale specification.
\par\emph{Why:} This makes the local clock a change of representation rather than a replacement for the physics, and fixes the quantities used to grade the coupled residual and energy closure.
\end{enumerate}
\subsection*{M1\_analytic\_clock}
\emph{Construct the exactly monotone local clock from a spatial coefficient network and analytic Gaussian error-function primitives.}
\begin{enumerate}[leftmargin=*,start=2]
  \item Create a `ClockBasisSpec` with the nondimensional time interval, fixed Gaussian centers \(`c_{j}`\), shared spacing `w`, floating-point type, `eps\_dtype`, strictly positive `kref`, the derived positive floor `kfloor`, and the choices `\{8,16,32\}`. 【author decision: lock \(`c_{j}`\), `w`, the floating-point type, and a strictly positive `kref` with units, source, and value before deriving \(`kfloor=sqrt(eps_{dtype})\times kref`\).】 Use a network \(`a_{psi}(x)`\) that reads position only and returns exactly \(`M_{tau}`\) coefficients;\allowbreak{} it must not read the pulse or any predicted field. 【author decision: declare the coefficient-network input normalization, layers, activations, and the scalar training-validation score used to select \(`M_{tau}`\);\allowbreak{} keep \(`M_{tau}=16`\) as the default and resolve exact ties toward the smaller candidate.】 Make each coefficient positive with softplus, evaluate the fixed Gaussian functions and their analytic error-function cumulative functions, and form `tau\_psi` exactly as E1. Obtain time and spatial derivatives from the analytic basis and automatic differentiation through the spatial network, never from a finite time integral. Return coefficients, basis values, the clock, all required clock derivatives, the chosen \(`M_{tau}`\), and the frozen basis specification in `ClockState`.
\noindent\emph{The clock is evaluated by analytic Gaussian primitives, never finite time quadrature. Fixed knots \(`c_{j}`\), spacing `w`, a declared positive reference rate `kref>0`, `kfloor=sqrt(eps\_dtype)*kref`, and a training-validation-only choice \(`M_{tau}\) in \{8,16,32\}` make `tau\_psi(x,0)=0` and \(`partial_{t}\) tau\_psi=kfloor+sum\_j \(softplus(a_{psi},j)b_{j}(t)>0`\) exactly for the implemented map.}
\par\noindent{\small\ttfamily [eq~1, raw LaTeX, not typeset] \textbackslash\{\}tau\_\textbackslash\{\}psi(\textbackslash\{\}mathbf\{x\},t)=k\_\{\textbackslash\{\}mathrm\{floor\}\}t+\textbackslash\{\}sum\_\{j=1\}\textasciicircum{}\{M\_\textbackslash\{\}tau\}\textbackslash\{\}operatorname\{softplus\}(a\_\{\textbackslash\{\}psi,j\}(\textbackslash\{\}mathbf\{x\}))B\_j(t),\textbackslash\{\}quad b\_j(t)=e\textasciicircum{}\{-((t-c\_j)/w)\textasciicircum{}2\},\textbackslash\{\}quad B\_j(t)=\textbackslash\{\}frac\{\textbackslash\{\}sqrt\{\textbackslash\{\}pi\}w\}\{2\}\textbackslash\{\}left[\textbackslash\{\}operatorname\{erf\}\textbackslash\{\}!\textbackslash\{\}left(\textbackslash\{\}frac\{t-c\_j\}\{w\}\textbackslash\{\}right)+\textbackslash\{\}operatorname\{erf\}\textbackslash\{\}!\textbackslash\{\}left(\textbackslash\{\}frac\{c\_j\}\{w\}\textbackslash\{\}right)\textbackslash\{\}right],\textbackslash\{\}quad k\_\{\textbackslash\{\}mathrm\{ref\}\}>0}
\par\emph{Why:} The analytic construction makes `tau\_psi(x,0) = 0` and the time derivative of `tau\_psi` strictly positive exactly, without finite quadrature. The forward graph is acyclic because the coefficient network never reads a predicted field.
\end{enumerate}
\subsection*{M2\_fields\_pullback\_residuals}
\emph{Predict the coupled fields, perform the complete pulse-aware derivative transformation, and evaluate the unchanged physics plus direct kinetics-clock residual.}
\begin{enumerate}[leftmargin=*,start=3]
  \item Store the pulse as a `PulseTrace`. On every smooth segment it must return the pulse value and its first two time derivatives;\allowbreak{} at each jump it must return separate left and right values with an \(`interface_{id}`\). Feed position, `tau\_psi`, and the current one-sided pulse value to one shared solution head \(`h_{theta}`\) with three outputs. 【author decision: declare the capacity-matched solution-head trunk, layers, activations, input normalization, and output parameterizations for voltage, temperature, and phase.】 Save those outputs as \(`V_{hat}`, `T_{hat}`\), and \(`xi_{hat}`\). Derive electric field from the physical spatial gradient of voltage, current density from the frozen conductivity, and Joule heating from current density dotted with electric field. Return the fields, these deterministic derived quantities, pulse record, clock reference, and physical-contract reference in `FieldState`.
\par\emph{Why:} One shared clock retains the electric-to-Joule-to-thermal-to-phase causal chain. The known one-sided pulse value enters only the solution head, so the clock graph remains acyclic.
  \item For each scalar field, build a `DerivativeBundle`. Inside a smooth pulse segment, use automatic differentiation to obtain derivatives with respect to position, `tau\_psi`, and every pulse component, including all mixed and second derivatives required by E3 and E4. Combine them with the first two pulse derivatives and the exact clock derivatives to form the complete first and second physical-time derivatives in E3. Form the complete \(`2\times 2`\) physical-space Hessian in E4, including both mixed position-clock terms, the clock-gradient outer product, and the clock Hessian. Keep the Hessian as a matrix;\allowbreak{} obtain a Laplacian only from its trace and a divergence only from the complete quantity being differentiated. At a pulse jump, evaluate separate left and right traces and never differentiate through the jump. Return all derivatives with their pulse-segment and side labels.
\noindent\emph{For each scalar `f=h\_theta(x,tau\_psi,p(t))`, the complete first- and second-order temporal pullback is evaluated on each open pulse-smooth interval. Pulse jumps are handled as interfaces and are never differentiated through.}
\par\noindent{\small\ttfamily [eq~2, raw LaTeX, not typeset] \textbackslash\{\}partial\_t f=f\_\textbackslash\{\}tau\textbackslash\{\}tau\_t+f\_p\textbackslash\{\}!\textbackslash\{\}cdot\textbackslash\{\}dot p,\textbackslash\{\}quad \textbackslash\{\}partial\_\{tt\}f=f\_\{\textbackslash\{\}tau\textbackslash\{\}tau\}\textbackslash\{\}tau\_t\textasciicircum{}2+f\_\textbackslash\{\}tau\textbackslash\{\}tau\_\{tt\}+2(f\_\{\textbackslash\{\}tau p\}\textbackslash\{\}!\textbackslash\{\}cdot\textbackslash\{\}dot p)\textbackslash\{\}tau\_t+f\_\{pp\}[\textbackslash\{\}dot p,\textbackslash\{\}dot p]+f\_p\textbackslash\{\}!\textbackslash\{\}cdot\textbackslash\{\}ddot p}
\noindent\emph{The full physical-space Hessian, including both mixed contractions, the clock-gradient outer product, and the clock Hessian, is constructed before taking a trace or divergence;\allowbreak{} a Hessian is never conflated with its Laplacian.}
\par\noindent{\small\ttfamily [eq~3, raw LaTeX, not typeset] D\_\{\textbackslash\{\}mathbf\{x\}\}\textasciicircum{}2f=f\_\{\textbackslash\{\}mathbf\{x\}\textbackslash\{\}mathbf\{x\}\}+f\_\{\textbackslash\{\}mathbf\{x\}\textbackslash\{\}tau\}\textbackslash\{\}otimes\textbackslash\{\}nabla\textbackslash\{\}tau\_\textbackslash\{\}psi+\textbackslash\{\}nabla\textbackslash\{\}tau\_\textbackslash\{\}psi\textbackslash\{\}otimes f\_\{\textbackslash\{\}tau\textbackslash\{\}mathbf\{x\}\}+f\_\{\textbackslash\{\}tau\textbackslash\{\}tau\}\textbackslash\{\}nabla\textbackslash\{\}tau\_\textbackslash\{\}psi\textbackslash\{\}otimes\textbackslash\{\}nabla\textbackslash\{\}tau\_\textbackslash\{\}psi+f\_\textbackslash\{\}tau D\_\{\textbackslash\{\}mathbf\{x\}\}\textasciicircum{}2\textbackslash\{\}tau\_\textbackslash\{\}psi,\textbackslash\{\}qquad \textbackslash\{\}Delta f=\textbackslash\{\}operatorname\{tr\}(D\_\{\textbackslash\{\}mathbf\{x\}\}\textasciicircum{}2f)}
\par\emph{Why:} A space-dependent time coordinate preserves the declared PDE only when every induced chain-rule term is retained with the correct tensor shape. Pulse jumps are treated as interfaces rather than differentiated through.
  \item On every interior point inside a smooth pulse segment, evaluate the three original physical residuals exactly as E5. Evaluate \(`K_{xi}`\) from all of its stated inputs: temperature, electric field, phase, the full physical phase gradient, and the physical phase Laplacian. Form \(`r_{tau}`\) directly from the exact clock derivative and the smooth positive kinetic magnitude in E2;\allowbreak{} do not introduce a log or centering step. Build the Physics-Informed Neural Network (PINN) objective in E5 with frozen residual scales and nonempty starting, boundary, material, and pulse-interface terms. At a pulse edge, pair left and right records at the same position and time, require temperature and phase continuity, and evaluate the electrical boundary condition separately on both sides so voltage may follow the imposed step. Form energy balance separately by integrating heat storage, conductive boundary flux, and Joule heating, and report it only as a diagnostic. Return every residual family, its weights and identifiers, and the separate energy record in `ResidualBatch`.
\noindent\emph{The sole auxiliary clock residual directly couples the exact clock derivative to the smooth positive phase-rate magnitude. The displayed normalized phase slope is conditional on vanishing phase and clock residuals;\allowbreak{} it is not an approximation bound.}
\par\noindent{\small\ttfamily [eq~4, raw LaTeX, not typeset] r\_\textbackslash\{\}tau=\textbackslash\{\}partial\_t\textbackslash\{\}tau\_\textbackslash\{\}psi-\textbackslash\{\}sqrt\{K\_\textbackslash\{\}xi(\textbackslash\{\}hat T,\textbackslash\{\}hat E,\textbackslash\{\}hat\textbackslash\{\}xi,\textbackslash\{\}nabla\textbackslash\{\}hat\textbackslash\{\}xi,\textbackslash\{\}Delta\textbackslash\{\}hat\textbackslash\{\}xi)\textasciicircum{}2+k\_\{\textbackslash\{\}mathrm\{floor\}\}\textasciicircum{}2\},\textbackslash\{\}qquad |\textbackslash\{\}partial\textbackslash\{\}hat\textbackslash\{\}xi/\textbackslash\{\}partial\textbackslash\{\}tau|=|K\_\textbackslash\{\}xi|/\textbackslash\{\}sqrt\{K\_\textbackslash\{\}xi\textasciicircum{}2+k\_\{\textbackslash\{\}mathrm\{floor\}\}\textasciicircum{}2\}<1\textbackslash\{\}quad\textbackslash\{\}text\{when \}r\_\textbackslash\{\}tau=0\textbackslash\{\}text\{ and \}\textbackslash\{\}mathcal\{R\}\_\textbackslash\{\}xi=0}
\noindent\emph{The original electric, Joule-heat, and phase equations are evaluated using the complete pulled-back derivatives, with frozen residual scales and nonempty initial, boundary, constitutive, and pulse-interface terms. The clock term is disabled when the identity control sets `tau=t`.}
\par\noindent{\small\ttfamily [eq~5, raw LaTeX, not typeset] \textbackslash\{\}mathbf\{R\}\_\{\textbackslash\{\}mathrm\{phys\}\}=\textbackslash\{\}left(\textbackslash\{\}nabla\textbackslash\{\}!\textbackslash\{\}cdot[\textbackslash\{\}sigma(\textbackslash\{\}hat T,\textbackslash\{\}hat\textbackslash\{\}xi)\textbackslash\{\}nabla\textbackslash\{\}hat V],\textbackslash\{\};\textbackslash\{\}rho c\_p\textbackslash\{\}partial\_t\textbackslash\{\}hat T-\textbackslash\{\}nabla\textbackslash\{\}!\textbackslash\{\}cdot[k(\textbackslash\{\}hat T,\textbackslash\{\}hat\textbackslash\{\}xi)\textbackslash\{\}nabla\textbackslash\{\}hat T]-\textbackslash\{\}sigma|\textbackslash\{\}nabla\textbackslash\{\}hat V|\textasciicircum{}2,\textbackslash\{\};\textbackslash\{\}partial\_t\textbackslash\{\}hat\textbackslash\{\}xi-K\_\textbackslash\{\}xi(\textbackslash\{\}hat T,\textbackslash\{\}hat E,\textbackslash\{\}hat\textbackslash\{\}xi,\textbackslash\{\}nabla\textbackslash\{\}hat\textbackslash\{\}xi,\textbackslash\{\}Delta\textbackslash\{\}hat\textbackslash\{\}xi)\textbackslash\{\}right),\textbackslash\{\}quad \textbackslash\{\}mathcal\{L\}=\textbackslash\{\}sum\_q\textbackslash\{\}|\textbackslash\{\}mathcal\{R\}\_q/s\_q\textbackslash\{\}|\_2\textasciicircum{}2+\textbackslash\{\}mathbf\{1\}\_\{\textbackslash\{\}mathrm\{clock\}\}\textbackslash\{\}|r\_\textbackslash\{\}tau/s\_\textbackslash\{\}tau\textbackslash\{\}|\_2\textasciicircum{}2+\textbackslash\{\}mathcal\{L\}\_\{\textbackslash\{\}mathrm\{IC/BC/interface\}\}}
\par\emph{Why:} This preserves the declared electro-thermal-phase model while making the clock specific to the kinetics. If both the phase residual and \(`r_{tau}`\) vanish, the magnitude of the phase slope with respect to clock time is below `1`;\allowbreak{} this is a conditional property, not an approximation guarantee.
\end{enumerate}
\subsection*{M3\_matched\_training}
\emph{Train the clock and fields under a capacity-, sample-, and budget-matched contract.}
\begin{enumerate}[leftmargin=*,start=6]
  \item Build a `CapacityPlan` that counts trainable clock and solution parameters and reduces only the solution-head allocation until the Kinetics-Clock Physics-Informed Neural Network (KC-PINN) total matches the declared reference allocation. Build one `TrainingContract` with nonempty point sets for the interior, starting state, boundaries, material conditions, and pulse interfaces;\allowbreak{} frozen residual scaling;\allowbreak{} joint clock-and-field updates from the first update;\allowbreak{} optimizer settings;\allowbreak{} a stopping rule;\allowbreak{} and one in-distribution checkpoint selector. 【author decision: lock the admissible solution-head size family and matching tolerance, all optimizer settings, update and stopping rules, and the single in-distribution checkpoint-selection criterion applied to every arm.】 Train with the unchanged step 5 (On every interior point inside \(a\dots )\) objective and save only checkpoints chosen by that rule. Attach parameter counts, point-set identifiers, scale identifiers, and the exact plan and contract versions to every checkpoint.
\par\emph{Why:} Reallocating capacity instead of adding it prevents parameter count or a different training contract from appearing to be a clock effect.
\end{enumerate}
\subsection*{M4\_validation}
\emph{Train matched baselines, run the identity intervention, lock all choices in distribution, and score untouched whole-case out-of-distribution cases.}
\begin{enumerate}[leftmargin=*,start=7]
  \item Describe every comparator in a frozen `BaselineSpec` and run it on the same training-validation case identifiers through the step 1 (Convert the \(existing\dots )\) and step 6 (Build a `CapacityPlan` that \(counts\dots )\) interfaces. 【author decision: provide a versioned implementation manifest for each comparator, including source, architecture or transform, trainable parameters, loss terms, and its exact mapping to the shared physical, sampling, and checkpoint-selection contracts.】 Include the same analytic monotone coordinate family with the complete pullback and unchanged phase residual containing \(`K_{xi}`\), while removing only the \(`K_{xi}`-driven\) clock quantities \(`k_{eff}`\) and \(`r_{tau}`\);\allowbreak{} an unmodified-coordinate multilayer-perceptron PINN;\allowbreak{} fixed analytic transforms;\allowbreak{} fixed-Fourier and adaptive-spectral PINNs;\allowbreak{} causal PINN and NeuSA;\allowbreak{} self-adaptive and loss-conditional weighting;\allowbreak{} and hard-constraint and AdamFLIP comparators. Also include a widened unmodified-coordinate arm with the same active trainable count as KC-PINN. Keep the original physical residual and pulse-interface evaluator in every arm and return checkpoints in the step 6 (Build a `CapacityPlan` that \(counts\dots )\) schema.
\par\emph{Why:} The candidate survives only if its kinetics-specific clock adds value beyond generic learned coordinates, analytic time warps, spectral capacity, causality, loss weighting, and hard constraints.
  \item Use the same model class with an `identity` clock setting. Create the same clock tensors and parameters, freeze them, replace their output with `tau\_psi=t`, and turn off the E5 clock term so \(`r_{tau}`\) is absent. Keep the solution head, named point sets, residual scales, derivative and pulse-interface code, optimizer settings, update and stopping rules, and checkpoint selector unchanged. Report both nominal and active trainable parameter counts. Dormant clock parameters give nominal parity, while the widened unmodified-coordinate arm in step 7 (Describe every comparator in \(a\dots )\) gives active-parameter parity. Return the identity checkpoint and a `ControlTrace` that records the switch, confirms the bypass, and lists shared contract identifiers.
\par\emph{Why:} This sole negative-control intervention tests whether downstream gains require the kinetics clock rather than dormant parameters or differently coded physics. The widened unmodified-coordinate arm in step 7 (Describe every comparator in \(a\dots )\) separately gives active-parameter-matched performance.
  \item Decode every checkpoint at requested physical times. For KC-PINN, evaluate the analytic `tau\_psi(x,t)` and then the solution head at position, clock time, and pulse value;\allowbreak{} do not invert the clock or interpolate along clock time. Create a `MetricSpec` for field alignment and error norm, switching event, phase-front extraction and Hausdorff distance, topology representation and discrepancy, energy-balance normalization, and coupled-residual normalization and aggregation. 【author decision: predeclare every event threshold, field norm, front level, Hausdorff and topology convention, energy normalization, and coupled-residual aggregation before finalizing the validation score.】 Apply the same physical-time metrics to KC-PINN, the identity arm, and every comparator on training-validation cases. Then freeze architecture, \(`M_{tau}`\), the clock basis, all other settings, thresholds, metric code, and stopping and checkpoint rules in `LockedScoringContract`.
\par\emph{Why:} Every quantity being estimated remains a physical-time functional, and no formal out-of-distribution result can influence model or evaluation choices.
  \item Load the versioned `LockedScoringContract` and an untouched `OODCaseManifest` with one complete held-out geometry, material combination, and pulse protocol per row. For each row, load the locked KC-PINN, identity, and comparator checkpoints without updating parameters, create only that case's physical inputs, and run the locked forward and decoding paths independently. Apply the exact step 9 (Decode every checkpoint \(at\dots )\) metrics to physical-time fields, switching, fronts, energy balance, and coupled residuals. Do not change any threshold, normalization, architecture, clock basis, comparator, or stopping and checkpoint rule after seeing a held-out result. Return one record per arm and whole case with case and contract identifiers and `evidence\_identity=synthetic\_numerical`;\allowbreak{} aggregate only by the locked rule and do not describe the result as experimental validation.
\par\emph{Why:} Whole-entity scoring after locking tests generalization without leakage and bounds any positive conclusion to synthetic numerical evidence for the declared VO2 model.
\end{enumerate}
\end{document}
